\chapter{Динамический family-order parameter и rank-one defect selector}

Постоянная антисимметричная форма на family triplet запрещена остаточной
группой, но динамическое поле
\begin{equation}
 \Sigma\in V=\chi_1\oplus\chi_2\oplus\chi_3
\end{equation}
обходит этот no-go. Поскольку \(\chi_1\chi_2\chi_3=1\), взаимодействие
\begin{equation}
 S_{\Sigma\chi}
 =\frac{ig}{2}\int_\gamma
 \epsilon_{ijk}\Sigma_k\chi_i\chi_j
\end{equation}
является \(G_{res}\)-инвариантным. Для \(\Sigma\ne0\) оператор
\begin{equation}
 A_{ij}(\Sigma)=g\epsilon_{ijk}\Sigma_k
\end{equation}
имеет rank (2) и kernel \(\mathbb R\Sigma\). Поэтому три локальные
Majorana-моды редуцируются к одной без ручной вставки rank-one projector.

\section{Исправление анализа потенциала}

Для изотропного тестового потенциала
\begin{equation}
 V(\Sigma)=m^2|\Sigma|^2+\frac u4|\Sigma|^4
 +v\Sigma_1\Sigma_2\Sigma_3,
 \qquad u>0,
\end{equation}
диагональные stationary points имеют
\begin{equation}
 \Sigma=s(\varepsilon_1,\varepsilon_2,\varepsilon_3),
 \qquad
 q=\varepsilon_1\varepsilon_2\varepsilon_3,
\end{equation}
где
\begin{equation}
 3us^2+qvs+2m^2=0.
\end{equation}
Следовательно, радиус не равен просто \(2|m^2|/u\) при \(v\ne0\).
Угловой Hessian положителен для \(qv<0\), поэтому знак \(v\) оставляет
четыре, а не восемь, энергетически предпочтительных ориентаций.

Более того, отрицательное bare \(m^2\) не является необходимым условием.
Для предпочтительной ветви первый порядок возникает при
\begin{equation}
 m^2=\frac{v^2}{27u},
\end{equation}
а ненулевой global minimum существует уже при
\begin{equation}
 m^2<\frac{v^2}{27u}.
\end{equation}
Таким образом, symmetry-allowed cubic term способен вызвать конденсацию даже
при положительной квадратичной массе. Главный динамический gate состоит не
только в доказательстве \(m^2<0\), а в выводе полного отношения
\((m^2,u,v)\) из parent action.

\section{Факторный quadratic selector}

Остаточная \(\mathbb Z_2\times\mathbb Z_2\) симметрия сама по себе разрешает
три независимых инварианта \(\Sigma_i^2\). Поэтому полностью изотропная
квадратичная масса является дополнительным предположением. Уже существующий
factor Laplacian даёт канонический кандидат
\begin{equation}
 M_\Sigma^2=m_0^2I+\kappa L(w_3,w_1),
\end{equation}
со спектром
\begin{equation}
 m_0^2+2\kappa w_3,
 \quad
 m_0^2+2\kappa w_1,
 \quad
 m_0^2+2\kappa(w_3+w_1).
\end{equation}
При \(w_1<w_3\) и \(\kappa>0\) первой теряет устойчивость уникальная
\(S^1\)-характерная компонента. Это геометрический mass selector, а не
ручной выбор оси. Открытым остаётся вывод коэффициента \(\kappa\) и его знака
из того же действия, которое создаёт defect.

\section{Вклад core Majorana pair}

Две моды, спаренные оператором \(A(\Sigma)\), имеют gap
\(g|\Sigma|\). Если разрешено минимизировать по обеим parity-ветвям
двухмодового core Hamiltonian, его zero-temperature ground-state energy
содержит
\begin{equation}
 \Delta E_f=-\frac12|g\Sigma|,
\end{equation}
а при конечном обратном температурном размере \(\beta\)
\begin{equation}
 \Delta F_f
 =-\frac1\beta\log\left[2\cosh
 \left(\frac{\beta g|\Sigma|}{2}\right)\right]
 =\mathrm{const}-\frac{\beta g^2}{8}|\Sigma|^2+O(|\Sigma|^4).
\end{equation}
Знак квадратичного вклада благоприятствует конденсации. Однако cusp
\(-|g\Sigma|/2\) не является автоматическим свойством полного defect.
При фиксированной глобальной fermion parity система может быть обязана
оставаться на одной ветви \(\pm g|\Sigma|/2\), а нечётное число локальных
Majorana operators требует bulk или удалённого parity completion. Кроме того,
при конечном \(\beta\) конденсация следует только после проверки знака полной
renormalized quadratic mass. Поэтому фермионный вклад создаёт реальный
механизм неустойчивости, но ещё не доказывает \(\Sigma\ne0\) при произвольном
bare \(m^2\).

Этот детерминант зависит только от \(|\Sigma|\) и не генерирует
\(\Sigma_1\Sigma_2\Sigma_3\); ориентационный cubic coefficient должен прийти
из Wilson/factor или другого bosonic boundary sector.

\section{Bulk--core compatibility gate}

Инвариант \(\epsilon_{ijk}\Sigma_k\chi_i\chi_j\) корректен как взаимодействие
уже спроецированных вещественных core Majorana operators. Его нельзя без
дополнительного вывода заменить bulk Yukawa-членом того же family-типа.
Для четырёхмерных Weyl/Majorana fields лоренц-скалярный bilinear
\(N_iN_j\) симметричен по family indices, поэтому contraction с
\(\epsilon_{ijk}\Sigma_k\) тождественно исчезает. Следовательно parent action
должен породить антисимметричный оператор либо после контролируемой
zero-mode projection, либо через family connection, derivative coupling или
иной bulk operator, чья core restriction равна \(A(\Sigma)\).

Это устраняет потенциальную circularity: локальные \(\chi_i\) нельзя считать
независимым фундаментальным входом до вывода самого defect spectrum.

\section{Точная граница ориентационной устойчивости}

Пусть эффективный радиальный cusp допустим и
\begin{equation}
 V_{\mathrm{eff}}=
 \sum_{i=1}^{3}\mu_i\Sigma_i^2+
 \frac u4|\Sigma|^4+v\Sigma_1\Sigma_2\Sigma_3-h|\Sigma|,
 \qquad h=\frac{|g|}{2}.
\end{equation}
Для осевого stationary point \(\Sigma=(0,s,0)\), \(s>0\), имеем
\begin{equation}
 2\mu_2s+us^3-h=0.
\end{equation}
Поперечный Hessian равен
\begin{equation}
 H_\perp=
 \begin{pmatrix}
 2(\mu_1-\mu_2) & vs\\
 vs & 2(\mu_3-\mu_2)
 \end{pmatrix}.
\end{equation}
Поэтому геометрически выбранная \(S^1\)-ось является локальным минимумом
точно при
\begin{equation}
 \mu_1>\mu_2,
 \qquad
 \mu_3>\mu_2,
 \qquad
 4(\mu_1-\mu_2)(\mu_3-\mu_2)>v^2s^2.
\end{equation}
Если последнее неравенство нарушено, cubic term уводит минимум в
многокомпонентную ветвь. Тем самым выражение ``анизотропия доминирует''
заменяется проверяемым условием.

\section{Статус rank-one утверждения}

В трёхмерном projected core subspace при любом \(\Sigma\ne0\) утверждение
точно: \(\operatorname{rank}A=2\) и \(\dim\ker A=1\). Для полного BdG operator
этого недостаточно. Нужен gap/Feshbach--Schur аргумент, исключающий
дополнительные bulk и core zero pairs и доказывающий, что coupling к высшим
модам не закрывает spectral gap. До такого доказательства rank-one является
exact projected result, но conditional full-spectrum result.

\begin{protocol}[Следующий parent-action gate]
Нужно вывести из одного действия \(m_0^2\), \(\kappa\), \(u\), разрешённые
анизотропные quartic coefficients, \(v\) и \(g\); получить антисимметричную
core coupling из допустимого bulk operator; зафиксировать global parity
completion; проверить поперечный Hessian; затем доказать full BdG gap.
Положительный исход допускает два механизма конденсации: отрицательный
fermionic mass shift или cubic-driven first-order transition. После
доказательства \(\Sigma\ne0\) rank-one kernel автоматически следует только
в projected core subspace.
\end{protocol}

\begin{nogocriterion}[Запрет ручного order parameter]
Если \(\Sigma\), знак \(\kappa\), cubic coefficient \(v\) или coupling
\(g\) назначаются после просмотра flavour/neutrino data, конструкция остаётся
новой EFT-гипотезой. Алгебраический selector при этом корректен, но физическая
ветвь не замкнута.
\end{nogocriterion}

Машинный аудит сохранён в
\texttt{s2t\_dynamic\_family\_order\_parameter\_results.json}.